% source: https://tex.stackexchange.com/a/749713 (question 749696, score 8, block 1)
\documentclass[12pt]{standalone}% compile with lualatex only
\usepackage[svgnames]{xcolor}
\usepackage{fouriernc}
\usepackage[3d]{luadraw}%https://github.com/pfradin/luadraw
%https://tex.stackexchange.com/questions/749696/drawing-a-3d-vector-field-with-vortices-and-perspective-axis-labels
% functions declarations
\begin{luacode}
function field(f,x1,x2,y1,y2,grid,length)  -- fonction mathématique, indépendante du graphique
-- calcule un champ de vecteurs dans le pavé [x1,x2]x[y1,y2]
-- f fonction de deux variables à valeurs dans R^2
-- grid = {nbx, nby} : nombre de vecteurs suivant x et suivant y
-- length = longueur d'un vecteur
    if grid == nil then grid = {25,25} end
    local deltax, deltay = (x2-x1)/(grid[1]-1), (y2-y1)/(grid[2]-1) -- pas suivant x et y
    if length == nil then length = math.min(deltax,deltay) end -- longueur par défaut
    local vectors = {} -- contiendra la liste des vecteurs
    local x, y, v = x1 
    for _ = 1, grid[1] do -- parcours suivant x
        y = y1
        for _ = 1, grid[2] do -- parcours suivant y
            v = f(x,y) -- on suppose que v est bien défini
            v = Z(v[1],v[2]) -- passage en complexe
            if not cpx.isNul(v) then
                v = v/cpx.abs(v)*length -- normalisation de v
                table.insert(vectors, {Z(x,y), Z(x,y)+v} ) -- on ajoute le vecteur
            end
            y = y+deltay
        end
        x = x+deltax
    end
    return vectors -- list of segments with complex numbers
end

function graph:Dvectorfield(f,args) -- add a method to graph
-- dessine un champ de vecteurs
-- f fonction de deux variables à valeurs dans R^2
-- args table à 4 champs :
-- { view={x1,x2,y1,y2}, grid={nbx,nby}, length=, draw_options=""}
    args = args or {}
    local view = args.view or {self:Xinf(),self:Xsup(),self:Yinf(),self:Ysup()} -- repère utilisateur par défaut
    local vectors = field(f,view[1],view[2],view[3],view[4],args.grid,args.length) -- calcul du champ
    self:Dpolyline(vectors,false,args.draw_options) -- le dessin (ligne polygonale non fermée)
end
\end{luacode}

\begin{document}
\begin{luadraw}{name=vectors_field}
local g = graph3d:new{window3d={-5,5,-5,5,-1,1},window={-8.5,8,-5,5},viewdir={30,60},size={12,6,0}}
g:Linejoin("round"); g:Labelsize("tiny"); g:Linewidth(4)
-- example with Volterra
local V = function(x,y) return {x-x*y,-y+x*y} end -- function: R^2 -> R^2
local H = function(t,Y) return V(Y[1],Y[2]) end -- associated differential equation Y'=H(t,Y)=V(Y)
g:Dboxaxes3d({grid=true,gridcolor="gray",fillcolor="LightGray",drawbox=true})
g:Savematrix()
g:Setmatrix( g:Proj3d({Origin,vecI,vecJ}) ) -- to draw in xy-plane instead of screen plane
g:Dvectorfield(V,{view={-4.5,4.5,-4.5,4.5},length=0.35,draw_options="line width=0.2pt,blue,-Stealth"})
g:Dodesolve(H,0,{2,3},{t={0,50},out={2,3},nbdots=250,draw_options="red,line width=0.8pt"})
g:Dodesolve(H,0,{-1,4},{t={0,4},out={2,3},nbdots=50,draw_options="green,line width=0.8pt"})
g:Restorematrix()
g:Show()
\end{luadraw}
\end{document}
